\section{A PRD szerepe a fejlesztési folyamatban}

A Product Requirements Document (PRD) a projekt előkészítésének egyik legfontosabb dokumentuma. Szerepe, hogy a termékötletet strukturált formába rendezze, és összekapcsolja a felhasználói igényeket, az üzleti célokat, a funkcionális követelményeket és a fejlesztési korlátokat. A PRD így nem pusztán leíró dokumentum, hanem a fejlesztési folyamat egyik kiindulópontja: segít meghatározni, hogy milyen problémát kell megoldani, kinek készül a termék, milyen funkciók szükségesek az első verzióhoz, és mely elemek kerülhetnek későbbi fejlesztési szakaszba.

A PRD elkészítése ideális esetben nem egyetlen személy feladata, hanem több szerepkör közös munkájának eredménye. A termékoldali szereplők a felhasználói igényeket és az üzleti célokat képviselik, a fejlesztők a technikai megvalósíthatóságot és a becsült ráfordítást vizsgálják, a tervezői és minőségbiztosítási szempontok pedig a használhatóságot, tesztelhetőséget és a követelmények egyértelműségét erősítik. Ez azért fontos, mert egy jól elkészített PRD nemcsak azt rögzíti, hogy mit kell megvalósítani, hanem azt is, hogy a csapat miért éppen ezeket a funkciókat választotta.

A WatchDB esetében a PRD különösen fontos szerepet tölt be, mert a kezdeti termékvízió szélesebb, mint az elsőként megvalósítható alkalmazás. A dokumentum alapján a projekt nemcsak személyes filmnaplóként értelmezhető, hanem hosszabb távon közösségi filmkövető platformként is. Ez a vízió több irányt tartalmaz: a felhasználó saját filmjeinek naplózását, a watchlist kezelését, értékelések rögzítését, valamint későbbi verziókban a barátok aktivitásának követését és a közösségi ajánlások megjelenítését.

A PRD egyik legfontosabb feladata ebben a helyzetben a kezdeti termék terjedelmének a kezelése. Egy teljes termékvízió gyakran több olyan funkciót tartalmaz, amelyek egyszerre történő megvalósítása túl nagy fejlesztési terhet jelentene. Ezért a dokumentum a funkciókat prioritás szerint csoportosítja: az MVP-hez szükséges funkciókra, a későbbi verziókban megvalósítható elemekre, valamint azokra a részekre, amelyek tudatosan kívül maradnak a projekt hatókörén. Ez a felosztás segít elkerülni, hogy az első verzió túl összetetté váljon, miközben megőrzi a termék hosszabb távú fejlesztési irányát.

A WatchDB PRD-je alapján az MVP központi célja, hogy a felhasználó gyorsan és egyszerűen tudjon filmeket keresni, megnézett filmeket naplózni, értékelést adni, valamint külön listában vezetni azokat a filmeket, amelyeket később szeretne megtekinteni. Ezek a funkciók alkotják azt a minimális terméket, amely már önmagában is értéket ad a felhasználónak, ez az MVP (Minimum Viable Product) lényege. A közösségi funkciók, például a publikus profilok, a követési rendszer vagy az aktivitási feed, a termékvízió fontos részei, de az első megvalósítható verzió szempontjából későbbi fejlesztési szakaszba sorolhatóak.

A PRD további előnye, hogy a fejlesztési döntéseket később visszakereshetővé és indokolhatóvá teszi. A célközönség, a felhasználói folyamatok, a nem funkcionális követelmények, a technológiai függőségek és a korlátok mind olyan tényezők, amelyek befolyásolják az alkalmazás tervezését. Például a külső filmadatok használata indokolja a TMDB API integrációját, míg az email-alapú bejelentkezés az MVP egyszerűsítését szolgálja.

Összességében a PRD a WatchDB fejlesztésében hidat képez a termékötlet és a technikai megvalósítás között. Segítségével a projekt nem elszigetelt funkciók gyűjteményeként jelenik meg, hanem olyan tudatosan szűkített fejlesztési feladatként, amely egy nagyobb termékvízió első, megvalósítható lépését képviseli és egy olyan irányt mutató dokumentum, amihez a csapat könnyen tud igazodni és utalni rá a fejlesztési folyamat során.
